\section{Introduction}\label{sec:intro}

Distributed LLM inference on edge clusters has moved from research
curiosity to viable deployment.  Petals~\cite{petals2023}
demonstrated that a collaborative swarm of consumer GPUs can serve
billion-parameter models, EdgeShard~\cite{edgeshard2024} added
latency-optimal layer placement via dynamic programming, and
Jupiter~\cite{jupiter2024} pushed throughput-optimal placement with
explicit Token-Between-Time (TBT) constraints.  Yet these systems
were built around two assumptions inherited from the data center
that the edge regime breaks.  First, \emph{homogeneity}: on our
8-worker Jetson fleet a Jetson AGX Orin in MAXN power mode is
16$\times$ faster per layer than a Jetson Orin Nano at decode-time
sequence lengths --- a finding that itself only emerged after we
fixed two silent profiler bugs (\S\ref{sec:eval:setup}).  Second,
\emph{reliability}: boards crash, SD cards fill up, networks
partition.  During the measurement campaign behind this paper we
encountered all three on multiple boards, and placement that
ignores them produces a system whose first worker failure aborts
the stream.

Existing systems treat layer placement~($\psi$) and recovery
routing~($R$) as separate problems.  EdgeShard's DP optimizes
$\psi$ assuming a redundant copy of every layer exists somewhere
on the cluster --- a comfortable assumption on a 16+\,GB peer but
unaffordable on a 4\,GB Nano.  Petals' greedy placement leaks
2.38$\times$ throughput on heterogeneous fleets~\cite{helix2024}
and its swarm-redundancy recovery model presupposes spare memory
the edge does not have.  Jupiter's throughput-optimal DP has no
recovery story at all.  Composing these systems naively does not
recover what each gives up: the systemic point we make in
\S\ref{sec:background} is that $R$ and $\psi$ have \emph{coupled}
feasibility regions on memory-tight edge.  Decoupling them gives
up feasibility or performance, never both.

\sys{} unifies $\psi$ and $R$ inside a single alternating dynamic
program whose backup-memory reservation folds back into the
placement feasibility check.  The same DP that decides which
worker hosts which transformer block also decides which worker
absorbs whose stage on failure, with one tunable cost knob $\alpha$
that lets the optimizer span EdgeShard's latency regime ($\alpha
{=}0$) and Jupiter's throughput regime ($\alpha{=}1$) under one
framework (\S\ref{sec:design:dp}).  Two runtime mechanisms make the
joint optimization viable on a real edge fleet: an
\emph{asynchronous mirror cache} that ships per-stage activations
back to the coordinator out-of-band, so a freshly-promoted backup
can rebuild its KV state without disrupting the surviving chain
(\S\ref{sec:design:mirror}); and \emph{asynchronous chain
forwarding} with a \texttt{ResultReady} reverse channel, which
breaks the synchronous nested-response unwind that we measured
silently serializes nominally-pipelined chains
(\S\ref{sec:design:async}).

We make the following contributions:

\begin{itemize}
  \item A formal $\psi$+$R$ alternating DP with polynomial-time
        convergence, generalizing EdgeShard's latency mode and
        Jupiter's throughput mode under one $\alpha$-tunable
        rank function (\S\ref{sec:design:dp}).
  \item An asynchronous mirror cache plus chain-aware failure
        attribution path that uses gRPC-trailer metadata to
        identify the truly-dead worker even when the heartbeat
        path has already substituted its backup
        (\S\ref{sec:design:mirror}).
  \item Asynchronous chain forwarding with per-request locking
        and a coordinator-side \texttt{ResultReady} reverse
        channel, yielding a chain-length-independent
        17--47\% throughput gain at $C{=}16$ over synchronous
        chains (\S\ref{sec:design:async}).
  \item Live measurement on an 8-worker Jetson fleet across
        \emph{28} operating points (two chain lengths $\times$
        four architecture variants $\times$ three--four
        concurrency levels) showing latency-optimal placement
        strictly dominates throughput-optimal placement at
        every point --- a structural property of the $\psi$+$R$
        joint optimization, not an artifact of any specific
        topology or runtime architecture
        (\S\ref{sec:eval:matrix}).  Under a mid-stream worker
        SIGKILL, \sys{} delivers all 60 expected tokens
        coherently while the placement-only baselines abort
        after the 17--20 tokens already in flight
        (\S\ref{sec:eval:recovery}).
\end{itemize}
